Throughout, $F=F(p,z,x)$ is smooth and
\[
F(Du,u,x)=0.
\]

\section{Complete Integrals, Envelopes}
\label{sec:complete-integrals-envelopes}

\subsection{Complete Integrals}
\label{sec:complete-integrals}

\begin{definition}[Complete integral]
\label{def:complete-integral}
\dcref{ch3:3.1.1-def-1}
\group{Nonlinear First-Order PDE}
A family $u=u(x,a)$, $a\in\R^n$, is a complete integral of
$F(Du,u,x)=0$ if every $u(\cdot,a)$ is a solution and
\[
\operatorname{rank}
\begin{pmatrix}D_a u\\D^2_{xa}u\end{pmatrix}=n.
\]
\end{definition}

\begin{remark}[Complete integrals depend on all $n$ parameters]
\label{rem:complete-integral-independence}
\uses{def:complete-integral}
\dcref{ch3:3.1.1-rem-1}
\group{Nonlinear First-Order PDE}

The rank condition excludes a nonzero direction $b\in\R^n$ for which
$D_a u\cdot b$ vanishes identically in $x$; differentiating such an identity
would put $b$ in the kernel of the displayed matrix.
\end{remark}

\begin{example}[Clairaut's equation]
\label{ex:clairaut-complete-integral}
\uses{def:complete-integral}
\dcref{ch3:3.1-ex-1}
\group{Nonlinear First-Order PDE}

For $x\cdot Du+f(Du)=u$, the functions
\[
u(x,a)=a\cdot x+f(a),\qquad a\in\R^n,
\]
form a complete integral.
\end{example}

\begin{example}[Eikonal equation]
\label{ex:eikonal-complete-integral}
\uses{def:complete-integral}
\dcref{ch3:3.1-ex-2}
\group{Nonlinear First-Order PDE}

For $|Du|=1$, a complete integral is
\[
u(x,a,b)=a\cdot x+b,qquad a\in S^{n-1},\ b\in\R.
\]
\end{example}

\begin{example}[Hamilton--Jacobi complete integral]
\label{ex:hamilton-jacobi-complete-integral}
\uses{def:complete-integral}
\dcref{ch3:3.1-ex-3}
\group{Hamilton-Jacobi Equations}

For $u_t+H(Du)=0$, the family
\[
u(x,t,a,b)=a\cdot x-tH(a)+b,qquad (a,b)\in\R^{n+1},
\]
is a complete integral.
\end{example}

\subsection{Envelopes and New Solutions}
\label{sec:envelopes-new-solutions}

\begin{definition}[Envelope]
\label{def:envelope}
\dcref{ch3:3.1.2-def-1}
\group{Nonlinear First-Order PDE}
Let $u=u(x,a)$ be a complete integral. If the equations
\[
D_a u(x,a)=0
\]
can be solved as $a=\phi(x)$, the function
\[
v(x):=u(x,\phi(x))
\]
is the envelope of the family.
\end{definition}

\begin{theorem}[Construction of new solutions]
\label{thm:envelope-solves-pde}
\uses{def:envelope}
\dcref{ch3:3.1-thm-1}
\group{Nonlinear First-Order PDE}

Every smooth envelope $v$ is a solution of $F(Dv,v,x)=0$.
\end{theorem}

\begin{proof}
\sketch The chain rule and $D_a u(x,\phi(x))=0$ give
$Dv=D_xu(x,\phi(x))$. Substitute this and $v=u(x,\phi(x))$ into the
equation satisfied by each member of the complete integral.
\end{proof}

\begin{remark}[Geometric interpretation of the envelope]
\label{rem:envelope-tangency}
\uses{thm:envelope-solves-pde}
\dcref{ch3:3.1.2-rem-1}
\group{Nonlinear First-Order PDE}

At each $x$, the graph of the envelope is tangent to the graph of the member
$u(\cdot,\phi(x))$.
\end{remark}

\begin{example}[Envelope of a complete integral for $u^2(1+|Du|^2)=1$]
\label{ex:envelope-singular-integral}
\uses{thm:envelope-solves-pde}
\dcref{ch3:3.1-ex-4}
\group{Nonlinear First-Order PDE}

The family
\[
u(x,a)=\pm\sqrt{1-|x-a|^2}
\]
solves $u^2(1+|Du|^2)=1$. The stationarity equation $D_a u=0$ gives
$a=x$, so its envelopes are the singular solutions $v\equiv\pm1$.
\end{example}

\begin{definition}[General integral]
\label{def:general-integral}
\uses{def:complete-integral, def:envelope}
\dcref{ch3:3.1.2-def-2}
\group{Nonlinear First-Order PDE}

Given a complete integral $u(x,a_1,\ldots,a_n)$ and a smooth function
$h:\R^{n-1}\to\R$, restrict to $a_n=h(a_1,\ldots,a_{n-1})$ and take the
envelope with respect to $a_1,\ldots,a_{n-1}$. The resulting family, with
arbitrary $h$, is the general integral associated with $u$.
\end{definition}

\begin{remark}[On the general integral]
\label{rem:general-integral-remarks}
\uses{def:general-integral}
\dcref{ch3:3.1.2-rem-2}
\group{Nonlinear First-Order PDE}

The construction yields solutions depending on one arbitrary function of
$n-1$ variables, but it need not contain every solution of a reducible or
factorized equation.
\end{remark}

\begin{example}[Alternative complete integral for the eikonal equation]
\label{ex:eikonal-alternative-complete-integral}
\uses{def:general-integral}
\dcref{ch3:3.1-ex-5}
\group{Nonlinear First-Order PDE}

In two dimensions,
\[
u(x,a_1,a_2)=x_1\cos a_1+x_2\sin a_1+a_2
\]
is a complete integral of $|Du|=1$. Taking $a_2=0$ and enveloping in $a_1$
gives $u=\pm|x|$ away from the origin.
\end{example}

\begin{example}[General integral for a Hamilton--Jacobi equation]
\label{ex:hamilton-jacobi-envelope-example}
\uses{def:general-integral, ex:hamilton-jacobi-complete-integral}
\dcref{ch3:3.1-ex-6}
\group{Hamilton-Jacobi Equations}

For $u_t+|Du|^2=0$, choose $b=0$ in the complete integral
$a\cdot x-t|a|^2+b$. Its envelope has $a=x/(2t)$ and hence
\[
u(x,t)=\frac{|x|^2}{4t}.
\]
\end{example}

\section{Characteristics}
\label{sec:characteristics}

\subsection{Derivation of the Characteristic ODE}
\label{sec:characteristic-ode-derivation}

Let $u$ solve $F(Du,u,x)=0$. Along a curve $x=x(s)$ set
\[
z(s)=u(x(s)),\qquad p(s)=Du(x(s)).
\]

\begin{theorem}[Structure of characteristic ODE]
\label{thm:characteristic-ode-structure}
\dcref{ch3:3.2-thm-1}
\group{Method of Characteristics}
The characteristic equations are
\[
\begin{cases}
\dot p=-D_xF(p,z,x)-F_z(p,z,x)p,\\
\dot z=D_pF(p,z,x)\cdot p,\\
\dot x=D_pF(p,z,x).
\end{cases}
\]
\end{theorem}

\begin{remark}[Closure of the characteristic system]
\label{rem:characteristic-ode-closure}
\uses{thm:characteristic-ode-structure}
\dcref{ch3:3.2.1-rem-1}
\group{Method of Characteristics}

The $2n+1$ equations form a closed first-order system for $(p,z,x)$.
\end{remark}

\subsection{Examples}
\label{sec:characteristics-examples}

For $b(x)\cdot Du+c(x)u=0$, the characteristic system reduces to
\[
\dot x=b(x),\qquad \dot z=-c(x)z.
\]

\begin{example}[Linear PDE on a quadrant]
\label{ex:linear-pde-quadrant-characteristics}
\uses{thm:characteristic-ode-structure}
\dcref{ch3:3.2-ex-1}
\group{Method of Characteristics}

The problem
\[
x_2u_{x_1}-x_1u_{x_2}=u\quad(x_1,x_2>0),
\qquad u(x_1,0)=g(x_1)
\]
has characteristic solution
\[
u(x_1,x_2)=g\!\left(\sqrt{x_1^2+x_2^2}\right)
\exp\!\left(\arctan\frac{x_2}{x_1}\right).
\]
\end{example}

For $b(x,u)\cdot Du+c(x,u)=0$, the reduced system is
\[
\dot x=b(x,z),\qquad \dot z=-c(x,z).
\]

\begin{example}[Semilinear PDE on a half-space]
\label{ex:semilinear-pde-halfspace-characteristics}
\uses{thm:characteristic-ode-structure}
\dcref{ch3:3.2-ex-2}
\group{Method of Characteristics}

For
\[
u_{x_1}+u_{x_2}=u^2\quad(x_2>0),\qquad u(x_1,0)=g(x_1),
\]
the characteristic formula is
\[
u(x_1,x_2)=\frac{g(x_1-x_2)}{1-x_2g(x_1-x_2)},
\]
where the denominator is nonzero.
\end{example}

\begin{example}[Fully nonlinear PDE on a half-space]
\label{ex:fully-nonlinear-halfspace-characteristics}
\uses{thm:characteristic-ode-structure}
\dcref{ch3:3.2-ex-3}
\group{Method of Characteristics}

The problem
\[
u_{x_1}u_{x_2}=u\quad(x_1>0),\qquad u(0,x_2)=x_2^2
\]
has the characteristic solution
\[
u(x_1,x_2)=\frac{(x_1+4x_2)^2}{16}.
\]
\end{example}

\subsection{Boundary Conditions}
\label{sec:characteristics-boundary-conditions}

After locally straightening the boundary, consider
\[
F(Du,u,x)=0\quad(x_n>0),\qquad u(x',0)=g(x').
\]

\begin{definition}[Compatibility conditions, admissible triple]
\label{def:compatibility-conditions-admissible-triple}
\dcref{ch3:3.2.3-def-1}
\group{Method of Characteristics}
A triple $(p^0,z^0,x^0)$ with $x_n^0=0$ is admissible if
\[
z^0=g(x^0),\qquad p_i^0=g_{x_i}(x^0)\ (i<n),\qquad
F(p^0,z^0,x^0)=0.
\]
\end{definition}

\begin{definition}[Noncharacteristic admissible triple]
\label{def:noncharacteristic-triple}
\uses{def:compatibility-conditions-admissible-triple}
\dcref{ch3:3.2.3-def-2}
\group{Method of Characteristics}

An admissible triple is noncharacteristic if
\[
F_{p_n}(p^0,z^0,x^0)\ne0.
\]
\end{definition}

\begin{lemma}[Noncharacteristic boundary conditions]
\label{lem:noncharacteristic-boundary-conditions}
\uses{def:noncharacteristic-triple}
\dcref{ch3:3.2-lem-1}
\group{Method of Characteristics}

Near a noncharacteristic admissible triple, the compatibility equation has a
unique smooth solution
\[
p_n^0=q(x^0,g(x^0),Dg(x^0)).
\]
\end{lemma}

\begin{proof}
\sketch Apply the implicit function theorem to
$F(g_{x_1}(x),\ldots,g_{x_{n-1}}(x),p_n,g(x),x)=0$; the derivative in
$p_n$ is nonzero by noncharacteristicity.
\end{proof}

\begin{remark}[Noncharacteristic condition for a general boundary]
\label{rem:noncharacteristic-condition-general}
\uses{def:noncharacteristic-triple}
\dcref{ch3:3.2.3-rem-1}
\group{Method of Characteristics}

For a smooth boundary with unit normal $\nu$, the invariant condition is
\[
D_pF(p^0,z^0,x^0)\cdot\nu(x^0)\ne0.
\]
\end{remark}

\subsection{Local Solution}
\label{sec:characteristics-local-solution}

For boundary coordinates $y\in\R^{n-1}$, solve the characteristic system
with initial data
\[
x(y,0)=(y,0),\quad z(y,0)=g(y),\quad
p(y,0)=(Dg(y),q(y,g(y),Dg(y))).
\]

\begin{lemma}[Local invertibility]
\label{lem:local-invertibility-characteristics}
\uses{def:noncharacteristic-triple}
\dcref{ch3:3.2-lem-2}
\group{Method of Characteristics}

The map $(y,s)\mapsto x(y,s)$ is locally invertible near $(y^0,0)$.
\end{lemma}

\begin{proof}
\sketch At $s=0$, the first $n-1$ columns of $D_{(y,s)}x$ are tangent to
the boundary, while the last is
$x_s=D_pF(p^0,z^0,x^0)$. Their determinant is nonzero precisely by the
noncharacteristic condition.
\end{proof}

Write the local inverse as $(y,s)=(y(x),s(x))$ and define
\[
u(x):=z(y(x),s(x)).
\]

\begin{theorem}[Local Existence Theorem]
\label{thm:local-existence-characteristics}
\uses{lem:local-invertibility-characteristics, def:compatibility-conditions-admissible-triple}
\dcref{ch3:3.2-thm-2}
\group{Method of Characteristics}

The function $u$ is smooth near $x^0$, satisfies
$F(Du,u,x)=0$, and agrees with $g$ on the boundary.
\end{theorem}

\begin{proof}
\sketch Along each characteristic, differentiation gives
\[
\frac d{ds}F(p,z,x)=D_pF\cdot(-D_xF-F_zp)
+F_z(D_pF\cdot p)+D_xF\cdot D_pF=0,
\]
so $F$ remains zero. The contact residuals
$z_s-p\cdot x_s$ and $z_{y_i}-p\cdot x_{y_i}$ vanish initially and satisfy
a homogeneous linear ODE system, hence vanish identically. The inverse chain
rule then gives $Du=p$, proving the PDE and the boundary condition.
\end{proof}

\subsection{Applications}
\label{sec:characteristics-applications}

\begin{example}[Flow patterns for a linear transport-type PDE]
\label{ex:linear-transport-characteristics-flow}
\uses{thm:local-existence-characteristics}
\dcref{ch3:3.2-ex-4}
\group{Method of Characteristics}

For $b(x)\cdot Du=0$, $u$ is constant along integral curves of $\dot x=b(x)$.
An attracting equilibrium obstructs a global continuous solution unless the
incoming boundary values coincide; a flow crossing the domain transports the
boundary values smoothly; and a characteristic equilibrium requires matching
values on the adjoining flow regions for continuity.
\end{example}

\begin{example}[Characteristics for conservation laws]
\label{ex:characteristics-conservation-laws}
\uses{thm:local-existence-characteristics}
\dcref{ch3:3.2-ex-5}
\group{Conservation Laws}

For
\[
u_t+\operatorname{div}F(u)=0,
\]
the projected characteristic from $x^0$ is
\[
x(t)=x^0+tF'(g(x^0)),\qquad u(x(t),t)=g(x^0).
\]
Crossing projected characteristics obstruct a global smooth solution.
\end{example}

\begin{remark}[Implicit formula for the conservation law]
\label{rem:implicit-formula-conservation-law}
\uses{ex:characteristics-conservation-laws}
\dcref{ch3:3.2.5-rem-1}
\group{Conservation Laws}

Before characteristic crossing, the solution satisfies
\[
u(x,t)=g(x-tF'(u(x,t))).
\]
The implicit function theorem applies while
\[
1+tDg(x-tF'(u))\cdot F''(u)\ne0.
\]
\end{remark}

\begin{example}[Characteristics for the Hamilton--Jacobi equation]
\label{ex:characteristics-hamilton-jacobi}
\uses{thm:characteristic-ode-structure}
\dcref{ch3:3.2-ex-6}
\group{Hamilton-Jacobi Equations}

For $u_t+H(Du,x)=0$, the characteristic system is
\[
\dot p=-D_xH(p,x),\qquad
\dot z=D_pH(p,x)\cdot p-H(p,x),\qquad
\dot x=D_pH(p,x).
\]
The first and third equations are Hamilton's equations.
\end{example}

\section{Introduction to Hamilton--Jacobi Equations}
\label{sec:hamilton-jacobi-intro}

We study
\[
u_t+H(Du)=0\quad\text{in }\R^n\times(0,\infty),\qquad
u(x,0)=g(x).
\]

\subsection{Calculus of Variations, Hamilton's ODE}
\label{sec:calculus-of-variations-hamiltons-ode}

For a smooth Lagrangian $L(q,x)$ and curves $w(0)=y$, $w(t)=x$, set
\[
I[w]=\int_0^tL(\dot w(s),w(s))\,ds.
\]

\begin{theorem}[Euler--Lagrange equations]
\label{thm:euler-lagrange-equations}
\dcref{ch3:3.3-thm-1}
\group{Calculus of Variations}
If $x(\cdot)$ minimizes $I$ among the admissible curves, then
\[
-\frac d{ds}D_qL(\dot x(s),x(s))+D_xL(\dot x(s),x(s))=0.
\]
\end{theorem}

\begin{proof}
\sketch For variations $x+\tau v$ with $v(0)=v(t)=0$, differentiating at
$\tau=0$ gives the first variation. Integration by parts and the fundamental
lemma of the calculus of variations yield the displayed system.
\end{proof}

\begin{remark}[Critical points versus minimizers]
\label{rem:critical-point-vs-minimizer}
\uses{thm:euler-lagrange-equations}
\dcref{ch3:3.3.1-rem-1}
\group{Calculus of Variations}

Every minimizer is a critical point of $I$, but a critical point need not be a
minimizer.
\end{remark}

\begin{example}[Newton's law as an Euler--Lagrange equation]
\label{ex:newtons-law-lagrangian}
\uses{thm:euler-lagrange-equations}
\dcref{ch3:3.3.1-ex-1}
\group{Calculus of Variations}

For $L(q,x)=\frac12m|q|^2-\phi(x)$, Euler--Lagrange gives
\[
m\ddot x=-D\phi(x)=:f(x).
\]
This is Newton's law (see \cite{FeynmanLeightonSands1966}, Chapter 19).
\end{example}

Set $p=D_qL(\dot x,x)$ and assume $p=D_qL(q,x)$ is uniquely solvable as
$q=\mathbf q(p,x)$.

\begin{definition}[Hamiltonian associated with a Lagrangian]
\label{def:hamiltonian-from-lagrangian}
\dcref{ch3:3.3.1-def-1}
\group{Calculus of Variations}
The Hamiltonian associated with $L$ is
\[
H(p,x)=p\cdot\mathbf q(p,x)-L(\mathbf q(p,x),x).
\]
\end{definition}

\begin{example}[Hamiltonian for Newton's law: total energy]
\label{ex:hamiltonian-newtons-law}
\uses{def:hamiltonian-from-lagrangian, ex:newtons-law-lagrangian}
\dcref{ch3:3.3.1-ex-2}
\group{Calculus of Variations}

For the Newtonian Lagrangian,
\[
H(p,x)=\frac1{2m}|p|^2+\phi(x),
\]
the total kinetic plus potential energy.
\end{example}

\begin{theorem}[Derivation of Hamilton's ODE]
\label{thm:hamiltons-ode-derivation}
\uses{thm:euler-lagrange-equations, def:hamiltonian-from-lagrangian}
\dcref{ch3:3.3-thm-2}
\group{Calculus of Variations}

The pair $(x,p)$ satisfies
\[
\dot x=D_pH(p,x),\qquad \dot p=-D_xH(p,x),
\]
and $H(p(s),x(s))$ is constant in $s$.
\end{theorem}

\begin{proof}
\sketch Differentiating $H=p\cdot q-L(q,x)$ and using $p=D_qL$ gives
$D_pH=q$ and $D_xH=-D_xL$. Insert $q=\dot x$ and the Euler--Lagrange
equation to obtain Hamilton's equations. Their contraction gives
$\frac d{ds}H=D_pH\cdot\dot p+D_xH\cdot\dot x=0$.
\end{proof}

\begin{remark}[Hamilton's ODE as a first-order system]
\label{rem:hamiltons-ode-system-structure}
\uses{thm:hamiltons-ode-derivation}
\dcref{ch3:3.3.1-rem-2}
\group{Calculus of Variations}

Hamilton's equations are a coupled system of $2n$ first-order ODE for
$(x,p)\in\R^{2n}$.
\end{remark}

\begin{remark}[Notational remark on classical mechanics]
\label{rem:arnold-mechanics-notation}
\dcref{ch3:3.3.1-rem-3}
\group{Calculus of Variations}
See \cite{Arnold1986}, Chapter 9, for the mechanics convention; the present
notation is chosen for PDE calculations.
\end{remark}

\subsection{Legendre Transform, Hopf--Lax Formula}
\label{sec:legendre-transform-hopf-lax}

Assume $L:\R^n\to\R$ is convex and superlinear:
\[
\lim_{|q|\to\infty}\frac{L(q)}{|q|}=+\infty.
\]

\begin{definition}[Legendre transform]
\label{def:legendre-transform}
\dcref{ch3:3.3.2-def-1}
\group{Calculus of Variations}
The Legendre transform of $L$ is
\[
L^*(p)=\sup_{q\in\R^n}\{p\cdot q-L(q)\}.
\]
\end{definition}

The maximizing relation $p=DL(q)$ identifies this transform with the
Hamiltonian; write $H=L^*$.

\begin{theorem}[Convex duality of Hamiltonian and Lagrangian]
\label{thm:convex-duality-hamiltonian-lagrangian}
\uses{def:legendre-transform}
\dcref{ch3:3.3-thm-3}
\group{Calculus of Variations}

The function $H=L^*$ is convex and superlinear, and
\[
L=H^*.
\]
\end{theorem}

\begin{proof}
\sketch A supremum of affine functions of $p$ is convex. Superlinearity of
$L$ yields superlinearity of $L^*$ by testing on $q=\lambda p/|p|$.
Fenchel's inequality gives $L\ge H^*$, while a supporting hyperplane to the
convex function $L$ at each $q$ gives the reverse inequality.
\end{proof}

\begin{remark}[Dual convex functions]
\label{rem:dual-convex-functions}
\uses{thm:convex-duality-hamiltonian-lagrangian}
\dcref{ch3:3.3.2-rem-1}
\group{Calculus of Variations}

$H$ and $L$ are dual convex functions: $H=L^*$ and $L=H^*$.
\end{remark}

\begin{theorem}[Hopf--Lax formula]
\label{thm:hopf-lax-formula-derivation}
\uses{thm:convex-duality-hamiltonian-lagrangian}
\dcref{ch3:3.3-thm-4}
\group{Hamilton-Jacobi Equations}

For $t>0$, the variational problem
\[
u(x,t)=\inf_{w(t)=x}\left\{\int_0^tL(\dot w(s))\,ds+g(w(0))\right\}
\]
equals
\[
u(x,t)=\min_{y\in\R^n}\left\{tL\left(\frac{x-y}{t}\right)+g(y)\right\}.
\]
\end{theorem}

\begin{proof}
\sketch The affine path from $y$ to $x$ gives the upper bound. Jensen's
inequality gives the reverse bound for every admissible path, and
superlinearity plus Lipschitz continuity of $g$ ensures the infimum is
attained.
\end{proof}

\begin{definition}[Hopf--Lax formula]
\label{def:hopf-lax-formula}
\uses{thm:hopf-lax-formula-derivation}
\dcref{ch3:3.3.2-def-2}
\group{Hamilton-Jacobi Equations}

The minimum expression in the preceding theorem is called the Hopf--Lax
formula.
\end{definition}

\begin{lemma}[A functional identity]
\label{lem:hopf-lax-functional-identity}
\uses{thm:hopf-lax-formula-derivation}
\dcref{ch3:3.3-lem-1}
\group{Hamilton-Jacobi Equations}

For $0\le s<t$,
\[
u(x,t)=\min_y\left\{(t-s)L\left(\frac{x-y}{t-s}\right)+u(y,s)\right\}.
\]
\end{lemma}

\begin{proof}
\sketch Substitute the Hopf--Lax representation of $u(y,s)$ and use
convexity of $L$ to combine the two action segments. The reverse inequality
uses a minimizing point $w$ at time $t$ and the intermediate point on the
affine segment from $w$ to $x$.
\end{proof}

\begin{lemma}[Lipschitz continuity]
\label{lem:hopf-lax-lipschitz-continuity}
\uses{thm:hopf-lax-formula-derivation}
\dcref{ch3:3.3-lem-2}
\group{Hamilton-Jacobi Equations}

If $g$ is Lipschitz, then $u$ is Lipschitz on $\R^n\times[0,\infty)$ and
$u(x,0)=g(x)$. In particular, for a constant $C$,
\[
|u(x,t)-g(x)|\le Ct,\qquad
|u(x,t)-u(\hat x,t)|\le \operatorname{Lip}(g)|x-\hat x|.
\]
\end{lemma}

\begin{proof}
\sketch Compare minimizers after translating $x$ to obtain the spatial
Lipschitz bound. Comparing with $y=x$ and using Fenchel duality bounds the
time increment by a constant times $|t-s|$; the same estimate at $t=0$ gives
the initial trace.
\end{proof}

\begin{theorem}[Solving the Hamilton--Jacobi equation]
\label{thm:hopf-lax-solves-hj-equation}
\uses{thm:hopf-lax-formula-derivation, lem:hopf-lax-functional-identity}
\dcref{ch3:3.3-thm-5}
\group{Hamilton-Jacobi Equations}

At every differentiability point $(x,t)$ with $t>0$, the Hopf--Lax function
satisfies
\[
u_t(x,t)+H(Du(x,t))=0.
\]
\end{theorem}

\begin{proof}
\sketch The functional identity with $s=t-h$ gives
$q\cdot Du+u_t\le L(q)$ for every $q$, hence
$u_t+H(Du)\le0$. A minimizing point $z$ in the Hopf--Lax formula and the
backward affine segment give the reverse inequality with
$q=(x-z)/t$. Take the supremum over $q$.
\end{proof}

\begin{theorem}[Hopf--Lax formula as solution]
\label{thm:hopf-lax-formula-as-solution}
\uses{thm:hopf-lax-solves-hj-equation, lem:hopf-lax-lipschitz-continuity}
\dcref{ch3:3.3-thm-6}
\group{Hamilton-Jacobi Equations}

The Hopf--Lax function is Lipschitz, differentiable almost everywhere, and
solves
\[
u_t+H(Du)=0\quad\text{a.e.},\qquad u(x,0)=g(x).
\]
\end{theorem}

\subsection{Weak Solutions, Uniqueness}
\label{sec:hj-weak-solutions-uniqueness}

\begin{example}[Nonuniqueness of Lipschitz a.e.\ solutions]
\label{ex:nonuniqueness-weak-solutions-hj}
\dcref{ch3:3.3.3-ex-1}
\group{Hamilton-Jacobi Equations}
For $u_t+|u_x|^2=0$ with $u(x,0)=0$, both $u_1\equiv0$ and
\[
u_2(x,t)=
\begin{cases}0,&|x|\ge t,\\x-t,&0\le x\le t,\\-x-t,&-t\le x\le0,
\end{cases}
\]
are Lipschitz a.e. solutions. Thus the a.e. notion alone is nonunique.
\end{example}

\begin{lemma}[Semiconcavity]
\label{lem:hopf-lax-semiconcavity-from-data}
\uses{def:semiconcave-function, thm:hopf-lax-formula-derivation}
\dcref{ch3:3.3-lem-3}
\group{Hamilton-Jacobi Equations}

If $g(x+z)-2g(x)+g(x-z)\le C|z|^2$, then its Hopf--Lax evolution obeys
\[
u(x+z,t)-2u(x,t)+u(x-z,t)\le C|z|^2.
\]
\end{lemma}

\begin{proof}
\sketch Use the same minimizing point $y$ in the formulas for $u(x+z,t)$,
$u(x,t)$, and $u(x-z,t)$. The Lagrangian terms cancel, leaving precisely the
second difference of $g$ at $y$.
\end{proof}

\begin{definition}[Semiconcave function]
\label{def:semiconcave-function}
\dcref{ch3:3.3.3-def-1}
\group{Hamilton-Jacobi Equations}
A function $g$ is semiconcave if
\[
g(x+z)-2g(x)+g(x-z)\le C|z|^2
\]
for some $C$ and all $x,z$. Equivalently, $g-C|x|^2/2$ is concave; bounded
$C^2$ Hessian is sufficient.
\end{definition}

\begin{remark}[Characterizations of semiconcavity]
\label{rem:semiconcave-characterization}
\uses{def:semiconcave-function}
\dcref{ch3:3.3.3-rem-1}
\group{Hamilton-Jacobi Equations}

The preceding second-difference characterization is equivalent to concavity
after subtracting a quadratic, and is implied by a globally bounded Hessian.
\end{remark}

\begin{definition}[Uniformly convex function]
\label{def:uniformly-convex-hamiltonian}
\dcref{ch3:3.3.3-def-2}
\group{Hamilton-Jacobi Equations}
A $C^2$ convex $H$ is uniformly convex with constant $\theta>0$ if
\[
\sum_{i,j}H_{p_ip_j}(p)\xi_i\xi_j\ge\theta|\xi|^2
\]
for all $p,\xi$.
\end{definition}

\begin{lemma}[Semiconcavity again]
\label{lem:hopf-lax-semiconcavity-uniform-convexity}
\uses{def:uniformly-convex-hamiltonian, thm:hopf-lax-formula-derivation}
\dcref{ch3:3.3-lem-4}
\group{Hamilton-Jacobi Equations}

If $H$ is uniformly convex with constant $\theta$, then
\[
u(x+z,t)-2u(x,t)+u(x-z,t)\le\frac{|z|^2}{\theta t}.
\]
\end{lemma}

\begin{proof}
\sketch Uniform convexity gives
$\frac12L(q_1)+\frac12L(q_2)\le L((q_1+q_2)/2)+|q_1-q_2|^2/(8\theta)$
by convex duality. Apply this estimate to the same minimizer in the three
Hopf--Lax expressions.
\end{proof}

\begin{definition}[Weak solution of the Hamilton--Jacobi initial-value problem]
\label{def:weak-solution-hamilton-jacobi}
\uses{def:semiconcave-function}
\dcref{ch3:3.3.3-def-3}
\group{Hamilton-Jacobi Equations}

A Lipschitz $u:\R^n\times[0,\infty)\to\R$ is a weak solution if
\begin{enumerate}
\item[(a)] $u(x,0)=g(x)$;
\item[(b)] $u_t+H(Du)=0$ almost everywhere for $t>0$;
\item[(c)] $u(x+z,t)-2u(x,t)+u(x-z,t)\le C(1+t^{-1})|z|^2$.
\end{enumerate}
\end{definition}

\begin{theorem}[Uniqueness of weak solutions]
\label{thm:uniqueness-weak-solutions-hj}
\uses{def:weak-solution-hamilton-jacobi}
\dcref{ch3:3.3-thm-7}
\group{Hamilton-Jacobi Equations}

If $H$ is $C^2$ and convex, superlinear, and $g$ is Lipschitz, there is at
most one weak solution.
\end{theorem}

\begin{proof}
\sketch For two solutions set $w=u-\tilde u$ and linearize
$H(Du)-H(D\tilde u)=b\cdot Dw$. Apply a nonnegative convex transform to $w$,
mollify in space-time, and use semiconcavity to obtain
$\operatorname{div}b_\varepsilon\le C(1+t^{-1})$. Integrating the resulting
transport inequality over backward cones gives a Gronwall estimate with zero
initial value. Letting the mollifier and transform parameters tend to zero
forces $w=0$.
\end{proof}

\begin{theorem}[Hopf--Lax formula as weak solution]
\label{thm:hopf-lax-formula-weak-solution}
\uses{lem:hopf-lax-semiconcavity-from-data, lem:hopf-lax-semiconcavity-uniform-convexity, thm:uniqueness-weak-solutions-hj}
\dcref{ch3:3.3-thm-8}
\group{Hamilton-Jacobi Equations}

If $H$ is $C^2$, convex and superlinear, $g$ is Lipschitz, and either $g$ is
semiconcave or $H$ uniformly convex, then the Hopf--Lax formula is the unique
weak solution.
\end{theorem}

\begin{example}[Explicit Hopf--Lax solution: $H(p) = \frac{1}{2}|p|^2$, $g = |x|$]
\label{ex:hopf-lax-explicit-example-i}
\uses{thm:hopf-lax-formula-weak-solution}
\dcref{ch3:3.3.3-ex-2}
\group{Hamilton-Jacobi Equations}

For $H(p)=|p|^2/2$, $g(x)=|x|$,
\[
u(x,t)=\min_y\left\{\frac{|x-y|^2}{2t}+|y|\right\}
=\begin{cases}|x|-t/2,&|x|\ge t,\\|x|^2/(2t),&|x|\le t.\end{cases}
\]
\end{example}

\begin{example}[Explicit Hopf--Lax solution: $H(p) = \frac{1}{2}|p|^2$, $g = -|x|$]
\label{ex:hopf-lax-explicit-example-ii}
\uses{thm:hopf-lax-formula-weak-solution}
\dcref{ch3:3.3.3-ex-3}
\group{Hamilton-Jacobi Equations}

For $g(x)=-|x|$, the same Hamiltonian gives
\[
u(x,t)=\min_y\left\{\frac{|x-y|^2}{2t}-|y|\right\}=-|x|-\frac t2.
\]
\end{example}

\section{Introduction to Conservation Laws}
\label{sec:conservation-laws-intro}

For the scalar conservation law
\[
u_t+F(u)_x=0\quad\text{in }\R\times(0,\infty),\qquad u(x,0)=g(x),
\]
smooth characteristics generally cross, so we use weak and entropy notions.

\subsection{Shocks, Entropy Condition}
\label{sec:shocks-entropy-condition}

For a compactly supported smooth test function $v$ on
$\R\times[0,\infty)$, the distributional initial-value identity is
\[
\int_0^\infty\!\int_\R(uv_t+F(u)v_x)\,dx\,dt+
\int_\R g(x)v(x,0)\,dx=0.
\]

\begin{definition}[Integral solution]
\label{def:integral-solution-conservation-law}
\dcref{ch3:3.4.1-def-1}
\group{Conservation Laws}
An $L^\infty$ function satisfying the preceding identity for every test
function is an integral solution of the conservation-law initial-value
problem.
\end{definition}

\begin{proposition}[Rankine--Hugoniot condition]
\label{prop:rankine-hugoniot-condition}
\uses{def:integral-solution-conservation-law}
\dcref{ch3:3.4.1-prop-1}
\group{Conservation Laws}

Suppose an integral solution is smooth on each side of a smooth discontinuity
curve $x=s(t)$, and that $u$ and its first derivatives extend uniformly
continuously to the curve from each side, with traces $u_l,u_r$. Then
\[
F(u_l)-F(u_r)=\dot s\,(u_l-u_r).
\]
Writing $[[u]]=u_l-u_r$, $[[F(u)]]=F(u_l)-F(u_r)$, and $\sigma=\dot s$,
this is $[[F(u)]]=\sigma[[u]]$.
\end{proposition}

\begin{proof}
\sketch Integrate the distributional identity separately on the two sides
of the curve. The interior terms vanish by the classical PDE; the boundary
fluxes combine to
$(F(u_l)-F(u_r))\nu^1+(u_l-u_r)\nu^2=0$. For the normal proportional to
$(1,-\dot s)$ this is the jump formula.
\end{proof}

\begin{example}[Shock waves for Burgers' equation]
\label{ex:burgers-shock-formation}
\uses{prop:rankine-hugoniot-condition}
\dcref{ch3:3.4-ex-1}
\group{Conservation Laws}

For $F(u)=u^2/2$ and
\[
g(x)=\begin{cases}1,&x\le0,\\1-x,&0\le x\le1,\\0,&x\ge1,\end{cases}
\]
characteristics give, for $0\le t<1$,
\[
u(x,t)=\begin{cases}
1,&x\le t,\\
(1-x)/(1-t),&t\le x\le1,\\
0,&x\ge1.
\end{cases}
\]
Thereafter the entropy continuation is the shock
\[
u=1\quad(x<s(t)),\qquad u=0\quad(x>s(t)),\qquad s(t)=\frac{1+t}{2},
\]
which satisfies Rankine--Hugoniot with speed $1/2$.
\end{example}

\begin{example}[Rarefaction waves and nonphysical shocks]
\label{ex:rarefaction-nonphysical-shock}
\uses{prop:rankine-hugoniot-condition, def:integral-solution-conservation-law}
\dcref{ch3:3.4-ex-2}
\group{Conservation Laws}

For Burgers data $g=0$ on $x<0$ and $g=1$ on $x>0$, both the discontinuous
profile with interface $x=t/2$ and the rarefaction
\[
u(x,t)=\begin{cases}0,&x<0,\\x/t,&0<x<t,\\1,&x>t\end{cases}
\]
are integral solutions. This nonuniqueness motivates an entropy condition.
\end{example}

For smooth solutions, a characteristic carries the value $g(x^0)$ along
$x=F'(g(x^0))t+x^0$. The physical requirement is that backward
characteristics do not cross a discontinuity.

\begin{definition}[Shock, entropy condition]
\label{def:shock-entropy-condition}
\uses{prop:rankine-hugoniot-condition}
\dcref{ch3:3.4.1-def-2}
\group{Conservation Laws}

At a discontinuity curve, a shock satisfies Rankine--Hugoniot and
\[
F'(u_l)>\sigma>F'(u_r).
\]
The inequalities are the entropy condition.
\end{definition}

\begin{remark}[Entropy condition under uniform convexity]
\label{rem:entropy-condition-uniform-convexity}
\uses{def:shock-entropy-condition}
\dcref{ch3:3.4.1-rem-2}
\group{Conservation Laws}

If $F''\ge\theta>0$, then $F'$ is strictly increasing and the entropy
inequalities are equivalent to $u_l>u_r$.
\end{remark}

\begin{example}[Shock and rarefaction combined for Burgers' equation]
\label{ex:burgers-triangular-shock-wave}
\uses{def:shock-entropy-condition, ex:burgers-shock-formation, ex:rarefaction-nonphysical-shock, prop:rankine-hugoniot-condition}
\dcref{ch3:3.4-ex-3}
\group{Conservation Laws}

For $g=0$ on $x<0$, $g=1$ on $0\le x\le1$, and $g=0$ on $x>1$, the
solution for $0\le t\le2$ is
\[
u=0\ (x<0),\quad u=x/t\ (0<x<t),\quad u=1\ (t<x<1+t/2),
\quad u=0\ (x>1+t/2).
\]
For $t\ge2$ the shock curve obeys
$\dot s=s/(2t)$, $s(2)=2$, hence $s(t)=\sqrt{2t}$, with $u=x/t$ to its
left and $u=0$ to its right.
\end{example}

\subsection{Lax--Oleinik Formula}
\label{sec:lax-oleinik-formula}

Assume $F$ is uniformly convex, $F(0)=0$, and set
\[
h(x)=\int_0^xg(y)\,dy,\qquad L=F^*.
\]
The Hopf--Lax function
\[
w(x,t)=\min_y\left\{tL\left(\frac{x-y}{t}\right)+h(y)\right\}
\]
satisfies $w_t+F(w_x)=0$ a.e. and $w(\cdot,0)=h$; formally $u=w_x$ solves
the conservation law.

Let $G=(F')^{-1}$.

\begin{theorem}[Lax--Oleinik formula]
\label{thm:lax-oleinik-formula}
\uses{def:legendre-transform, thm:hopf-lax-formula-derivation}
\dcref{ch3:3.4-thm-1}
\group{Conservation Laws}

For each $t>0$, except at most countably many $x$, the minimizer $y(x,t)$
in the Hopf--Lax formula is unique; $x\mapsto y(x,t)$ is nondecreasing; and
\[
u(x,t)=G\left(\frac{x-y(x,t)}{t}\right)
\]
for almost every $(x,t)$.
\end{theorem}

\begin{proof}
\sketch Since $L'(q)=G(q)$ and $L''>0$, strict convexity gives existence of
minimizers and the monotonicity of the smallest minimizer by comparing two
spatial points. At points of differentiability, differentiate the minimized
expression; stationarity in $y$ cancels the $y_x$ terms and yields
$u=L'((x-y)/t)=G((x-y)/t)$.
\end{proof}

\begin{definition}[Lax--Oleinik formula]
\label{def:lax-oleinik-formula}
\uses{thm:lax-oleinik-formula}
\dcref{ch3:3.4.2-def-1}
\group{Conservation Laws}

The preceding representation is the Lax--Oleinik formula.
\end{definition}

\begin{theorem}[Lax--Oleinik formula as integral solution]
\label{thm:lax-oleinik-integral-solution}
\uses{thm:lax-oleinik-formula, def:integral-solution-conservation-law}
\dcref{ch3:3.4-thm-2}
\group{Conservation Laws}

The Lax--Oleinik function is an integral solution of the conservation-law
initial-value problem.
\end{theorem}

\begin{proof}
\sketch Apply the a.e. Hamilton--Jacobi equation for $w$ to a test function
derivative $v_x$, integrate by parts in $x$ and $t$, and use $w(\cdot,0)=h$
and $h_x=g$. Since $u=w_x$ a.e., the resulting identity is the integral
solution identity.
\end{proof}

\subsection{Weak Solutions, Uniqueness}
\label{sec:conservation-laws-entropy-uniqueness}

\begin{lemma}[A one-sided jump estimate]
\label{lem:lax-oleinik-entropy-estimate}
\uses{thm:lax-oleinik-formula}
\dcref{ch3:3.4-lem-1}
\group{Conservation Laws}

There is $C$ such that the Lax--Oleinik solution satisfies
\[
u(x+z,t)-u(x,t)\le\frac Ct z\qquad(z>0).
\]
\end{lemma}

\begin{proof}
\sketch Restrict minimizers to a bounded velocity interval, so $G$ is
Lipschitz there. Monotonicity of $y(x,t)$ and $G$ then give
\[
G((x-y(x,t))/t)\ge G((x+z-y(x+z,t))/t)-\operatorname{Lip}(G)z/t.
\]
This is the claimed estimate.
\end{proof}

\begin{definition}[Entropy condition]
\label{def:entropy-condition-inequality}
\uses{lem:lax-oleinik-entropy-estimate}
\dcref{ch3:3.4.3-def-1}
\group{Conservation Laws}

The preceding one-sided estimate is called the entropy condition.
\end{definition}

\begin{remark}[Consequences of the entropy condition]
\label{rem:entropy-condition-consequences}
\uses{def:entropy-condition-inequality}
\dcref{ch3:3.4.3-rem-1}
\group{Conservation Laws}

For each $t>0$, $u(x,t)-Cx/t$ is nonincreasing. Thus left and right traces
exist everywhere and at a discontinuity $u_l\ge u_r$, agreeing with the shock
entropy inequality under uniform convexity.
\end{remark}

\begin{definition}[Entropy solution]
\label{def:entropy-solution-conservation-law}
\uses{def:entropy-condition-inequality}
\dcref{ch3:3.4.3-def-2}
\group{Conservation Laws}

An $L^\infty$ function is an entropy solution if it satisfies the integral
identity and, for some $C\ge0$,
\[
u(x+z,t)-u(x,t)\le C(1+t^{-1})z
\]
for almost every $x,t$ and every $z>0$.
\end{definition}

\begin{theorem}[Uniqueness of entropy solutions]
\label{thm:uniqueness-entropy-solutions}
\uses{def:entropy-solution-conservation-law}
\dcref{ch3:3.4-thm-3}
\group{Conservation Laws}

If $F$ is smooth and convex, entropy solutions are unique up to a null set.
\end{theorem}

\begin{proof}
\sketch For two entropy solutions set $w=u-\bar u$ and
$b=\int_0^1F'(ru+(1-r)\bar u)\,dr$, so the difference identity is
$\int w(v_t+bv_x)=0$. Mollify both solutions, solve the backward linear
transport equation with velocity $b_\varepsilon$ for a compactly supported
source, and use the one-sided entropy bounds to control
$b_{\varepsilon,x}\le C/t$ and the total variation near $t=0$. First let
$\varepsilon\to0$, then the initial time cutoff tend to zero; duality gives
$\int w\psi=0$ for every test source.
\end{proof}

\subsection{Riemann's Problem}
\label{sec:riemann-problem}

For the piecewise constant data
\[
g(x)=\begin{cases}u_l,&x<0,\\u_r,&x>0,\end{cases}
\]
assume $F$ is $C^2$ and uniformly convex and write $G=(F')^{-1}$.

\begin{theorem}[Solution of Riemann's problem]
\label{thm:riemann-problem-solution}
\uses{def:shock-entropy-condition, thm:uniqueness-entropy-solutions, prop:rankine-hugoniot-condition}
\dcref{ch3:3.4-thm-4}
\group{Conservation Laws}

If $u_l>u_r$, the unique entropy solution is the shock
\[
u(x,t)=\begin{cases}u_l,&x/t<\sigma,\\u_r,&x/t>\sigma,\end{cases}
\qquad \sigma=\frac{F(u_l)-F(u_r)}{u_l-u_r}.
\]
If $u_l<u_r$, the unique entropy solution is the rarefaction
\[
u(x,t)=\begin{cases}
u_l,&x/t\le F'(u_l),\\
G(x/t),&F'(u_l)\le x/t\le F'(u_r),\\
u_r,&x/t\ge F'(u_r).
\end{cases}
\]
\end{theorem}

\begin{proof}
\sketch In the shock case, Rankine--Hugoniot gives $\sigma$ and strict
convexity places $F'(u_r)<\sigma<F'(u_l)$, so entropy holds. In the rarefaction
case, substituting $u=v(x/t)$ into the PDE gives
$v'(x/t)[F'(v)-x/t]/t=0$, hence $v=G(x/t)$ between the two constant states;
continuity and the one-sided estimate verify the entropy condition.
\end{proof}

\begin{remark}[Interpretation of the two cases]
\label{rem:riemann-problem-shock-vs-rarefaction}
\uses{thm:riemann-problem-solution}
\dcref{ch3:3.4.4-rem-1}
\group{Conservation Laws}

The first case is a constant-speed shock and the second a centered
rarefaction. Both coincide with the Lax--Oleinik construction by uniqueness.
\end{remark}

\subsection{Long Time Behavior}
\label{sec:conservation-laws-long-time-behavior}

Assume $F$ is smooth and uniformly convex, $F(0)=0$, and $g$ is bounded and
integrable.

\begin{theorem}[Asymptotics in $L^\infty$-norm]
\label{thm:asymptotics-sup-norm}
\uses{thm:lax-oleinik-formula}
\dcref{ch3:3.4-thm-5}
\group{Conservation Laws}

There is $C$ such that
\[
|u(x,t)|\le Ct^{-1/2}\qquad(x\in\R,t>0).
\]
\end{theorem}

\begin{proof}
\sketch Put $\sigma=F'(0)$, so $G(\sigma)=L(\sigma)=L'(\sigma)=0$.
Uniform convexity gives
$tL((x-y)/t)\ge\theta|x-y-\sigma t|^2/t$. Boundedness of
$h=\int_0^xg$ forces every minimizer to satisfy
$|(x-y(x,t))/t-\sigma|\le Ct^{-1/2}$. The Lax--Oleinik formula and the
Lipschitz continuity of $G$ yield the result.
\end{proof}

For the N-wave asymptotics, assume in addition that $g$ has compact support.

\begin{definition}[N-wave]
\label{def:n-wave}
\dcref{ch3:3.4.5-def-1}
\group{Conservation Laws}
For $p,q\ge0$, $d>0$, and velocity $\sigma$, define
\[
N(x,t)=\begin{cases}
\dfrac1d\left(\dfrac{x}{t}-\sigma\right),
&-(pdt)^{1/2}<x-\sigma t<(qdt)^{1/2},\\
0,&\text{otherwise}.
\end{cases}
\]
\end{definition}

Let $\sigma=F'(0)$, $d=F''(0)$, and
\[
p=-2\min_y\int_{-\infty}^y g(x)\,dx,\qquad
q=2\max_y\int_y^\infty g(x)\,dx.
\]
Then $p,q\ge0$ and $G'(\sigma)=1/d$.

\begin{theorem}[Asymptotics in $L^1$-norm]
\label{thm:asymptotics-l1-norm-n-wave}
\uses{thm:asymptotics-sup-norm, def:n-wave}
\dcref{ch3:3.4-thm-6}
\group{Conservation Laws}

If $p,q>0$, there is $C$ such that
\[
\int_\R|u(x,t)-N(x,t)|\,dx\le Ct^{-1/2}.
\]
\end{theorem}

\begin{proof}
\sketch Taylor expansion of $G$ at $\sigma$ and the minimizer bound give
\[
\left|u(x,t)-\frac1d\frac{(x-\sigma t)-y(x,t)}{t}\right|\le Ct^{-1}.
\]
Compact support of $g$ makes $h$ constant outside a bounded interval and
identifies its two minima through $p,q$. Comparing the Hopf--Lax functional
with the endpoint constants localizes the support of $u$ to the two N-wave
edges up to $O(1)$, while monotonicity of $y(x,t)$ controls the interior.
Together with the $L^\infty$ estimate this gives the stated $L^1$ rate.
\end{proof}

\begin{example}[Burgers' triangular shock wave as an $N$-wave]
\label{ex:burgers-example-n-wave-specialization}
\uses{def:n-wave, ex:burgers-triangular-shock-wave, thm:asymptotics-l1-norm-n-wave}
\dcref{ch3:3.4.5-ex-1}
\group{Conservation Laws}

For the triangular Burgers data, $p=0$, $q=2$, $\sigma=0$, and $d=1$.
Thus
\[
N(x,t)=\begin{cases}x/t,&0<x<(2t)^{1/2},\\0,&\text{otherwise},\end{cases}
\]
and the entropy solution equals this N-wave for $t\ge2$.
\end{example}

\section{References}
\label{sec:ch3-references}

\noindent\cite{CourantHilbert1962,Caratheodory1982,Garabedian1964,John1982,Chester1971,Sneddon1957,Zwillinger1984,PLLions1982,Rund1973,Benton1977,Benton1977,Lax1973,Smoller1983,DiPerna1983,Whitham1974}
